\begin{abstract}
Reinforcement learning (RL) is increasingly applied to traffic engineering
(TE), but whether TE decisions are sequential enough to need it is rarely
measured. We study an incremental MPLS-TE problem in which routing state
persists, one label-switched path may move per five-minute interval, moves
are costly, a hold-down follows each move and an action mask enforces
failures and protected-class capacity---formally a Markov decision process.
On a flow-level 18-router backbone, a masked contextual bandit that predicts
only the immediate reward of each action beats MaskablePPO on all
\TestBanditRoots{} training roots (paired return difference \Gap{}, 95\,\% CI
[\GapLo, \GapHi]) and matches a non-learning per-interval min-max-utilization
MILP controller. The cause is the planning horizon, not the algorithm: PPO with
$\gamma=0$ matches the bandit, while PPO with $\gamma=0.995$ oscillates at
hold-down expiry and loses its network gains to reconfiguration costs, and a
bootstrapped value learner does not improve on the bandit. A clairvoyant
lookahead diagnostic explains why: with traffic held fixed, the
first-interval best action is the 24-interval best action in
\SeqFrozenAgreeHTwentyFour\,\% of states, against \SeqLiveAgreeHTwentyFour\,\%
under the true future, so almost all non-myopic value is anticipation of
exogenous change that the controller cannot observe. Delaying the effect of a
move by one interval reverses the ranking (PPO $-$ bandit $=\DelayPPOVs$
[\DelayPPOVsLo, \DelayPPOVsHi]). We also show that a prior internal study's
bandit result reproduces on different hardware while its PPO result does not.
Sequential RL for TE should be justified by measured temporal coupling, not by
the formal presence of state.
\end{abstract}
